%% preview.tex — 図の確認用スタンドアロン文書。
%% ルートディレクトリから platex+dvipdfmx でコンパイルする:
%%   platex -output-directory=analysis -halt-on-error analysis/preview.tex
%%   dvipdfmx -o analysis/preview.pdf analysis/preview.dvi
%% latexmk が analysis/ を作業ディレクトリにする場合にも対応。
\documentclass[dvipdfmx,a4paper]{jsarticle}
\usepackage{amsmath,amssymb,bm}
\makeatletter
\def\input@path{{analysis/}}
\makeatother
% latexmk がソース側へ移動して実行する場合と、ルートから直接 platex する場合の両方に対応。
\IfFileExists{analysis/data/transport_m3.dat}{%
  \newcommand{\figdata}{analysis/data}%
}{%
  \newcommand{\figdata}{data}%
}
%% figstyle.tex
%% 卒論用 pgfplots 図の共通設定。
%% 本文 (1029307812_特別研究報告書2026.tex) のプリアンブルで
%%   %% figstyle.tex
%% 卒論用 pgfplots 図の共通設定。
%% 本文 (1029307812_特別研究報告書2026.tex) のプリアンブルで
%%   %% figstyle.tex
%% 卒論用 pgfplots 図の共通設定。
%% 本文 (1029307812_特別研究報告書2026.tex) のプリアンブルで
%%   \input{analysis/figstyle}
%% を呼び、各図は \input{analysis/fig_xxx} で取り込む。
%%
%% ※ documentclass のオプションに [dvipdfmx] が付いているので、
%%   tikz/pgfplots は自動で dvipdfmx ドライバを使う（platex + dvipdfmx 構成）。

\usepackage{tikz}
\usepackage{pgfplots}
\usepackage{pgfplotstable}
\pgfplotsset{compat=1.18}
\usepgfplotslibrary{groupplots}
\usetikzlibrary{pgfplots.groupplots}

%% データファイルの場所（本文からの相対パス）。必要なら上書き可。
\providecommand{\figdata}{analysis/data}

%% 図が多い節で float が後ろの節へ流れ込まないように配置を緩める。
\renewcommand{\topfraction}{0.9}
\renewcommand{\bottomfraction}{0.9}
\renewcommand{\textfraction}{0.08}
\renewcommand{\floatpagefraction}{0.75}
\setcounter{topnumber}{2}
\setcounter{bottomnumber}{1}
\setcounter{totalnumber}{4}

%% '#' 始まりの行をコメントとして無視させる（gnuplot 風 .dat 用）。
\pgfplotstableset{col sep=space, comment chars={\#}}

%% beta = 0.5,1,2,3 を表す色（色覚多様性に配慮した4色）
\definecolor{betaA}{HTML}{0072B2} % beta=0.5  青
\definecolor{betaB}{HTML}{009E73} % beta=1    緑
\definecolor{betaC}{HTML}{E69F00} % beta=2    橙
\definecolor{betaD}{HTML}{D55E00} % beta=3    朱

%% 図の基本スタイル
\pgfplotsset{
  thesisaxis/.style={
    width=7.2cm, height=5.6cm,
    scale only axis,
    tick align=outside,
    tick label style={font=\footnotesize},
    label style={font=\small},
    legend style={font=\footnotesize, draw=none, fill opacity=0.85,
      text opacity=1, row sep=1pt, inner sep=2pt},
    every axis plot/.append style={line width=0.9pt, mark size=1.6pt},
    grid=both,
    grid style={line width=0.2pt, draw=gray!18},
    major grid style={line width=0.3pt, draw=gray!30},
  },
  % beta 4本のカーブを一括設定する系列スタイル
  betaplotA/.style={color=betaA, mark=*,        mark options={fill=betaA}},
  betaplotB/.style={color=betaB, mark=square*,  mark options={fill=betaB}},
  betaplotC/.style={color=betaC, mark=triangle*,mark options={fill=betaC}},
  betaplotD/.style={color=betaD, mark=diamond*, mark options={fill=betaD}},
}

%% delta 軸（対数）の共通設定
\pgfplotsset{
  deltaxaxis/.style={
    xmode=log, log basis x=10,
    xtick={0.3333,1,2,3,5,10},
    xticklabels={$\tfrac13$,$1$,$2$,$3$,$5$,$10$},
    xlabel={$\frac{|\Delta \alpha|E}{p}$},
    xmin=0.28, xmax=12,
  },
}

%% beta 凡例（4系列に共通で使う）
\newcommand{\betalegend}{\legend{$\beta pE=0.5$,$\beta pE=1$,$\beta pE=2$,$\beta pE=3$}}

%% 4本の beta カーブ（誤差棒つき）を一括描画。
%%   #1 データファイル, #2 x列名, #3 値の接頭辞, #4 誤差の接頭辞
\newcommand{\plotbetaserr}[4]{%
  \addplot[betaplotA, error bars/.cd, y dir=both, y explicit]
    table[x=#2, y=#3_b05, y error=#4_b05]{#1};
  \addplot[betaplotB, error bars/.cd, y dir=both, y explicit]
    table[x=#2, y=#3_b1,  y error=#4_b1]{#1};
  \addplot[betaplotC, error bars/.cd, y dir=both, y explicit]
    table[x=#2, y=#3_b2,  y error=#4_b2]{#1};
  \addplot[betaplotD, error bars/.cd, y dir=both, y explicit]
    table[x=#2, y=#3_b3,  y error=#4_b3]{#1};
}

%% 4本の beta カーブ（誤差棒なし）。#1 ファイル, #2 x列, #3 値の接頭辞
\newcommand{\plotbetas}[3]{%
  \addplot[betaplotA] table[x=#2, y=#3_b05]{#1};
  \addplot[betaplotB] table[x=#2, y=#3_b1]{#1};
  \addplot[betaplotC] table[x=#2, y=#3_b2]{#1};
  \addplot[betaplotD] table[x=#2, y=#3_b3]{#1};
}

%% を呼び、各図は \input{analysis/fig_xxx} で取り込む。
%%
%% ※ documentclass のオプションに [dvipdfmx] が付いているので、
%%   tikz/pgfplots は自動で dvipdfmx ドライバを使う（platex + dvipdfmx 構成）。

\usepackage{tikz}
\usepackage{pgfplots}
\usepackage{pgfplotstable}
\pgfplotsset{compat=1.18}
\usepgfplotslibrary{groupplots}
\usetikzlibrary{pgfplots.groupplots}

%% データファイルの場所（本文からの相対パス）。必要なら上書き可。
\providecommand{\figdata}{analysis/data}

%% 図が多い節で float が後ろの節へ流れ込まないように配置を緩める。
\renewcommand{\topfraction}{0.9}
\renewcommand{\bottomfraction}{0.9}
\renewcommand{\textfraction}{0.08}
\renewcommand{\floatpagefraction}{0.75}
\setcounter{topnumber}{2}
\setcounter{bottomnumber}{1}
\setcounter{totalnumber}{4}

%% '#' 始まりの行をコメントとして無視させる（gnuplot 風 .dat 用）。
\pgfplotstableset{col sep=space, comment chars={\#}}

%% beta = 0.5,1,2,3 を表す色（色覚多様性に配慮した4色）
\definecolor{betaA}{HTML}{0072B2} % beta=0.5  青
\definecolor{betaB}{HTML}{009E73} % beta=1    緑
\definecolor{betaC}{HTML}{E69F00} % beta=2    橙
\definecolor{betaD}{HTML}{D55E00} % beta=3    朱

%% 図の基本スタイル
\pgfplotsset{
  thesisaxis/.style={
    width=7.2cm, height=5.6cm,
    scale only axis,
    tick align=outside,
    tick label style={font=\footnotesize},
    label style={font=\small},
    legend style={font=\footnotesize, draw=none, fill opacity=0.85,
      text opacity=1, row sep=1pt, inner sep=2pt},
    every axis plot/.append style={line width=0.9pt, mark size=1.6pt},
    grid=both,
    grid style={line width=0.2pt, draw=gray!18},
    major grid style={line width=0.3pt, draw=gray!30},
  },
  % beta 4本のカーブを一括設定する系列スタイル
  betaplotA/.style={color=betaA, mark=*,        mark options={fill=betaA}},
  betaplotB/.style={color=betaB, mark=square*,  mark options={fill=betaB}},
  betaplotC/.style={color=betaC, mark=triangle*,mark options={fill=betaC}},
  betaplotD/.style={color=betaD, mark=diamond*, mark options={fill=betaD}},
}

%% delta 軸（対数）の共通設定
\pgfplotsset{
  deltaxaxis/.style={
    xmode=log, log basis x=10,
    xtick={0.3333,1,2,3,5,10},
    xticklabels={$\tfrac13$,$1$,$2$,$3$,$5$,$10$},
    xlabel={$\frac{|\Delta \alpha|E}{p}$},
    xmin=0.28, xmax=12,
  },
}

%% beta 凡例（4系列に共通で使う）
\newcommand{\betalegend}{\legend{$\beta pE=0.5$,$\beta pE=1$,$\beta pE=2$,$\beta pE=3$}}

%% 4本の beta カーブ（誤差棒つき）を一括描画。
%%   #1 データファイル, #2 x列名, #3 値の接頭辞, #4 誤差の接頭辞
\newcommand{\plotbetaserr}[4]{%
  \addplot[betaplotA, error bars/.cd, y dir=both, y explicit]
    table[x=#2, y=#3_b05, y error=#4_b05]{#1};
  \addplot[betaplotB, error bars/.cd, y dir=both, y explicit]
    table[x=#2, y=#3_b1,  y error=#4_b1]{#1};
  \addplot[betaplotC, error bars/.cd, y dir=both, y explicit]
    table[x=#2, y=#3_b2,  y error=#4_b2]{#1};
  \addplot[betaplotD, error bars/.cd, y dir=both, y explicit]
    table[x=#2, y=#3_b3,  y error=#4_b3]{#1};
}

%% 4本の beta カーブ（誤差棒なし）。#1 ファイル, #2 x列, #3 値の接頭辞
\newcommand{\plotbetas}[3]{%
  \addplot[betaplotA] table[x=#2, y=#3_b05]{#1};
  \addplot[betaplotB] table[x=#2, y=#3_b1]{#1};
  \addplot[betaplotC] table[x=#2, y=#3_b2]{#1};
  \addplot[betaplotD] table[x=#2, y=#3_b3]{#1};
}

%% を呼び、各図は \input{analysis/fig_xxx} で取り込む。
%%
%% ※ documentclass のオプションに [dvipdfmx] が付いているので、
%%   tikz/pgfplots は自動で dvipdfmx ドライバを使う（platex + dvipdfmx 構成）。

\usepackage{tikz}
\usepackage{pgfplots}
\usepackage{pgfplotstable}
\pgfplotsset{compat=1.18}
\usepgfplotslibrary{groupplots}
\usetikzlibrary{pgfplots.groupplots}

%% データファイルの場所（本文からの相対パス）。必要なら上書き可。
\providecommand{\figdata}{analysis/data}

%% 図が多い節で float が後ろの節へ流れ込まないように配置を緩める。
\renewcommand{\topfraction}{0.9}
\renewcommand{\bottomfraction}{0.9}
\renewcommand{\textfraction}{0.08}
\renewcommand{\floatpagefraction}{0.75}
\setcounter{topnumber}{2}
\setcounter{bottomnumber}{1}
\setcounter{totalnumber}{4}

%% '#' 始まりの行をコメントとして無視させる（gnuplot 風 .dat 用）。
\pgfplotstableset{col sep=space, comment chars={\#}}

%% beta = 0.5,1,2,3 を表す色（色覚多様性に配慮した4色）
\definecolor{betaA}{HTML}{0072B2} % beta=0.5  青
\definecolor{betaB}{HTML}{009E73} % beta=1    緑
\definecolor{betaC}{HTML}{E69F00} % beta=2    橙
\definecolor{betaD}{HTML}{D55E00} % beta=3    朱

%% 図の基本スタイル
\pgfplotsset{
  thesisaxis/.style={
    width=7.2cm, height=5.6cm,
    scale only axis,
    tick align=outside,
    tick label style={font=\footnotesize},
    label style={font=\small},
    legend style={font=\footnotesize, draw=none, fill opacity=0.85,
      text opacity=1, row sep=1pt, inner sep=2pt},
    every axis plot/.append style={line width=0.9pt, mark size=1.6pt},
    grid=both,
    grid style={line width=0.2pt, draw=gray!18},
    major grid style={line width=0.3pt, draw=gray!30},
  },
  % beta 4本のカーブを一括設定する系列スタイル
  betaplotA/.style={color=betaA, mark=*,        mark options={fill=betaA}},
  betaplotB/.style={color=betaB, mark=square*,  mark options={fill=betaB}},
  betaplotC/.style={color=betaC, mark=triangle*,mark options={fill=betaC}},
  betaplotD/.style={color=betaD, mark=diamond*, mark options={fill=betaD}},
}

%% delta 軸（対数）の共通設定
\pgfplotsset{
  deltaxaxis/.style={
    xmode=log, log basis x=10,
    xtick={0.3333,1,2,3,5,10},
    xticklabels={$\tfrac13$,$1$,$2$,$3$,$5$,$10$},
    xlabel={$\frac{|\Delta \alpha|E}{p}$},
    xmin=0.28, xmax=12,
  },
}

%% beta 凡例（4系列に共通で使う）
\newcommand{\betalegend}{\legend{$\beta pE=0.5$,$\beta pE=1$,$\beta pE=2$,$\beta pE=3$}}

%% 4本の beta カーブ（誤差棒つき）を一括描画。
%%   #1 データファイル, #2 x列名, #3 値の接頭辞, #4 誤差の接頭辞
\newcommand{\plotbetaserr}[4]{%
  \addplot[betaplotA, error bars/.cd, y dir=both, y explicit]
    table[x=#2, y=#3_b05, y error=#4_b05]{#1};
  \addplot[betaplotB, error bars/.cd, y dir=both, y explicit]
    table[x=#2, y=#3_b1,  y error=#4_b1]{#1};
  \addplot[betaplotC, error bars/.cd, y dir=both, y explicit]
    table[x=#2, y=#3_b2,  y error=#4_b2]{#1};
  \addplot[betaplotD, error bars/.cd, y dir=both, y explicit]
    table[x=#2, y=#3_b3,  y error=#4_b3]{#1};
}

%% 4本の beta カーブ（誤差棒なし）。#1 ファイル, #2 x列, #3 値の接頭辞
\newcommand{\plotbetas}[3]{%
  \addplot[betaplotA] table[x=#2, y=#3_b05]{#1};
  \addplot[betaplotB] table[x=#2, y=#3_b1]{#1};
  \addplot[betaplotC] table[x=#2, y=#3_b2]{#1};
  \addplot[betaplotD] table[x=#2, y=#3_b3]{#1};
}


\newcommand{\figexplain}[1]{%
  \par\vspace{0.8em}%
  \noindent\textbf{この図から説明できること：}#1\par
}

\begin{document}
\begin{center}\large 卒論用 図プレビュー\end{center}

\section*{Fig.~1: $v$ と $D_{\mathrm{eff}}$ の $\frac{|\Delta \alpha|E}{p}$ 依存性}
\begin{center}%% fig_transport.tex
%% 主結果: 平均速度 v=L/T1 と有効拡散係数 D_eff の |Delta alpha|E/p 依存性。
%% 列 = 棒長 (m=3, m=6)、行 = (v, D_eff)。誤差棒はブートストラップ標準誤差。
\begin{tikzpicture}
  \begin{groupplot}[
      group style={group size=2 by 2, horizontal sep=2.0cm,
        vertical sep=1.7cm},
      thesisaxis, deltaxaxis,
      width=6.15cm, height=5.1cm,
    ]
    % --- 左上: v, m=3 ---
    \nextgroupplot[ylabel={$v=L/T_1$}, title={$m=3$ (棒長 $\approx$ 最狭窄部)},
      ymin=0.25, ymax=1.05, legend pos=north west]
      \plotbetaserr{\figdata/transport_m3.dat}{delta}{v}{sev}
      \draw[densely dashed, gray] ({axis cs:1,0}|-{rel axis cs:0,0})
        -- ({axis cs:1,0}|-{rel axis cs:0,1});
      \betalegend
    % --- 右上: v, m=6 ---
    \nextgroupplot[ylabel={$v=L/T_1$}, title={$m=6$ (棒長 $\approx 2\times$ 最狭窄部)},
      ymin=0.1, ymax=0.65]
      \plotbetaserr{\figdata/transport_m6.dat}{delta}{v}{sev}
      \draw[densely dashed, gray] ({axis cs:2,0}|-{rel axis cs:0,0})
        -- ({axis cs:2,0}|-{rel axis cs:0,1});
    % --- 左下: D_eff, m=3 ---
    \nextgroupplot[ylabel={$D_{\mathrm{eff}}$}, ymin=0.1, ymax=0.42]
      \plotbetaserr{\figdata/transport_m3.dat}{delta}{D}{seD}
    % --- 右下: D_eff, m=6 ---
    \nextgroupplot[ylabel={$D_{\mathrm{eff}}$}, ymin=0.05, ymax=0.25]
      \plotbetaserr{\figdata/transport_m6.dat}{delta}{D}{seD}
  \end{groupplot}
\end{tikzpicture}
\end{center}
\figexplain{誘起双極子の寄与を表す $\frac{|\Delta \alpha|E}{p}$ が大きくなるほど、棒がチャネル軸方向へ傾き、狭窄部を通りやすくなるため、平均速度 $v$ と有効拡散係数 $D_{\mathrm{eff}}$ が増加することを説明できる。棒が長い $m=6$ ではこの変化がより急で、幾何的な通りにくさが強く効いている。}
\clearpage

\section*{Fig.~2: もっと流路の幅が大きい場所での $\langle|\sin\phi|\rangle$}
\begin{center}%% fig_orientation_sin.tex
%% 機構の定量化: チャネルの広い断面 (x~0.20) における占有重み付き <|sin phi|>
%% （横方向への張り出し l|sin phi| の指標）の |Delta alpha|E/p 依存性。
%% 安定配向 sin(phi1)=p/(|Delta alpha|E) が与える理論曲線を重ねる。
\begin{tikzpicture}
  \begin{groupplot}[
      group style={group size=2 by 1, horizontal sep=2.0cm},
      thesisaxis, deltaxaxis,
      width=6.25cm,
      ymin=0, ymax=0.95,
      ylabel={$\langle|\sin\phi|\rangle$ (広い断面)},
    ]
    \nextgroupplot[title={$m=3$}, legend pos=south west]
      % 理論: 安定配向の |sin phi1| = min(1, p/(|Delta alpha|E))
      \addplot[black, densely dashed, line width=1pt, domain=0.28:1, samples=2,
        forget plot]{1};
      \addplot[black, densely dashed, line width=1pt, domain=1:11, samples=60]
        {1/x};
      \plotbetas{\figdata/transport_m3.dat}{delta}{sinw}
      \legend{{$\min\!\left(1,\frac{p}{|\Delta \alpha|E}\right)$},{$\beta pE=0.5$},{$\beta pE=1$},%
        {$\beta pE=2$},{$\beta pE=3$}}
    \nextgroupplot[title={$m=6$}, ylabel={$\langle|\sin\phi|\rangle$ (広い断面)}]
      \addplot[black, densely dashed, line width=1pt, domain=0.28:1, samples=2,
        forget plot]{1};
      \addplot[black, densely dashed, line width=1pt, domain=1:11, samples=60,
        forget plot]{1/x};
      \plotbetas{\figdata/transport_m6.dat}{delta}{sinw}
  \end{groupplot}
\end{tikzpicture}
\end{center}
\figexplain{この図はチャネルの最も幅が広い場所における $\sin \phi$ の滞在時間による重み付き平均 $\langle|\sin\phi|\rangle$ を、熱ノイズがない場合の平衡解 $\frac{p}{|\Delta \alpha| E}$ と共に表した図である。$\sin \phi$ は棒を $y$ 軸方向へ射映した大きさの割合である。この値が $\frac{|\Delta \alpha|E}{p}$ の増加とともに小さくなることから、電場が棒を $x$ 軸方向に平行な向きに傾けていることがわかる。}
\clearpage

\section*{Fig.~3: 角度の分布 $P(\phi)$}
\begin{center}%% fig_pphi.tex
%% 機構の可視化: 広い断面での配向分布 P(phi)。
%% delta が小さいと phi=90deg（電場方向・最大張り出し）に単峰、
%% delta が大きいと phi1=arcsin(1/delta) 付近（軸方向）へ二峰化して移動する。
\begin{tikzpicture}
  \begin{groupplot}[
      group style={group size=2 by 1, horizontal sep=1.4cm},
      thesisaxis,
      xmin=0, xmax=360,
      xtick={0,90,180,270,360},
      xticklabels={$0$,$\tfrac{\pi}{2}$,$\pi$,$\tfrac{3\pi}{2}$,$2\pi$},
      xlabel={$\phi$}, ylabel={$P(\phi)$},
      ymin=0, ymax=0.62,
      cycle list={
        {betaA, line width=1pt},
        {betaB, line width=1pt},
        {betaC, line width=1pt},
        {betaD, line width=1pt}},
    ]
    \nextgroupplot[title={$m=3,\ \beta pE=3$}]
      % phi=90deg（小 delta の安定配向）の目印
      \draw[gray, densely dotted] (axis cs:90,0) -- (axis cs:90,0.62);
      \addplot table[x=phi_deg, y=d0333333]{\figdata/pphi_wide_m3_b3.dat};
      \addplot table[x=phi_deg, y=d1]{\figdata/pphi_wide_m3_b3.dat};
      \addplot table[x=phi_deg, y=d3]{\figdata/pphi_wide_m3_b3.dat};
      \addplot table[x=phi_deg, y=d10]{\figdata/pphi_wide_m3_b3.dat};
      \legend{$\delta=\tfrac13$,$\delta=1$,$\delta=3$,$\delta=10$}
    \nextgroupplot[title={$m=6,\ \beta pE=3$}, ylabel={}]
      \draw[gray, densely dotted] (axis cs:90,0) -- (axis cs:90,0.62);
      \addplot table[x=phi_deg, y=d0333333]{\figdata/pphi_wide_m6_b3.dat};
      \addplot table[x=phi_deg, y=d1]{\figdata/pphi_wide_m6_b3.dat};
      \addplot table[x=phi_deg, y=d3]{\figdata/pphi_wide_m6_b3.dat};
      \addplot table[x=phi_deg, y=d10]{\figdata/pphi_wide_m6_b3.dat};
  \end{groupplot}
\end{tikzpicture}
\end{center}
\figexplain{$\frac{|\Delta \alpha|E}{p}$ が小さいと棒は電場方向 $\phi=\pi/2$ に集中し、値が大きいと $\phi=0,\pi$ に近い軸方向へ分布が移る。この図から、輸送量の増加が角度分布の変化によって起きていることを説明できる。}
\clearpage

\section*{Fig.~4 角度と速度の対応}
\begin{center}%% fig_collapse.tex
%% 配向 → 輸送の対応: 横軸に広い断面での <|sin phi|>（張り出しの指標）、
%% 縦軸に v=L/T1。全 48 combo を散布。単調減少の帯に乗れば
%% 「配向（張り出し）が輸送を支配している」ことの直接的証拠になる。
\begin{tikzpicture}
  \begin{axis}[
      thesisaxis,
      width=8.0cm, height=6.2cm,
      xlabel={$\langle|\sin\phi|\rangle$ (広い断面)},
      ylabel={$v=L/T_1$},
      xmin=0.1, xmax=0.85, ymin=0.1, ymax=1.05,
      legend pos=north east,
    ]
    \addplot[only marks, mark=*, mark size=1.8pt, color=betaA,
      mark options={fill=betaA, fill opacity=0.7}]
      table[x=sin_wide, y=v]{\figdata/collapse_m3.dat};
    \addplot[only marks, mark=square*, mark size=1.8pt, color=betaC,
      mark options={fill=betaC, fill opacity=0.7}]
      table[x=sin_wide, y=v]{\figdata/collapse_m6.dat};
    \legend{$m=3$,$m=6$}
  \end{axis}
\end{tikzpicture}
\end{center}
\figexplain{$y$ 軸への射映成分 $\langle|\sin\phi|\rangle$ が小さい条件ほど平均速度 $v$ が大きい。同じ張り出しでも $m=6$ は $m=3$ より遅く、角度に加えて棒長そのものも輸送抵抗を大きくすることがわかる。}
\clearpage

\section*{Fig.~5 分離の動作マップ $v\!\left(\frac{|\Delta \alpha|E}{p},\ \beta pE\right)$}
\begin{center}%% fig_heatmap.tex
%% 分離の「動作マップ」: パラメータ平面 (|Delta alpha|E/p, beta pE) 上での
%% 平均速度 v=L/T1 を上から見たヒートマップで表す。同じ動作点でも
%% m（棒長）で v が大きく異なれば、速度（移動度）の差として
%% 電気的・幾何的性質を分離できることを示す。
\begin{tikzpicture}
  \begin{groupplot}[
      group style={group size=2 by 1, horizontal sep=4.1cm},
      width=4.25cm, height=4.35cm, scale only axis,
      view={0}{90},                      % 真上から見た 2D ヒートマップ
      tick align=outside,
      tick label style={font=\footnotesize}, label style={font=\small},
      title style={font=\small},
      xlabel={$\frac{|\Delta \alpha|E}{p}$},
      xtick={-1.585,0,1,1.585,2.322,3.322},
      xticklabels={$\tfrac13$,$1$,$2$,$3$,$5$,$10$},
      ylabel={$\beta pE$}, ytick={0.5,1,2,3},
      xmin=-1.585, xmax=3.322, ymin=0.5, ymax=3,
      enlargelimits=false,
      colormap/viridis,
      colorbar, colorbar style={font=\footnotesize, width=0.22cm,
        ylabel={$v=L/T_1$}},
    ]
    \nextgroupplot[title={$m=3$}, point meta min=0.25, point meta max=1.0]
      \addplot3[surf, shader=interp, mesh/cols=6]
        table[x=log2delta, y=beta, z=v]{\figdata/heat_m3.dat};
    \nextgroupplot[title={$m=6$}, point meta min=0.13, point meta max=0.62]
      \addplot3[surf, shader=interp, mesh/cols=6]
        table[x=log2delta, y=beta, z=v]{\figdata/heat_m6.dat};
  \end{groupplot}
\end{tikzpicture}
\end{center}
\figexplain{電場の強さや粒子の電気的性質を表すパラメータ平面上で、どの条件なら速く進むかを読み取れる。同じ動作点でも棒長によって $v$ が異なるため、移動度差を利用した分離条件を選べることを説明できる。この図はわかりづらいから論文に入れなくてもいいかも。}
\clearpage

\section*{Fig.~6 $(x,\phi)$ 占有マップ $P(\phi\,|\,x)$}
\begin{center}%% fig_xphi.tex
%% (x, phi) 占有マップ: 各位置 x（1周期に畳んだもの）での配向条件分布 P(phi|x)。
%% m=6, beta pE=3 で、delta=1/3（配向が電場方向に揃う）と delta=10（軸方向へ傾く）を比較。
%% 破線は最狭窄部 x~0.81 と最広部 x~0.19。点線は phi=pi/2（電場方向）。
\begin{tikzpicture}
  \begin{groupplot}[
      group style={group size=1 by 2, vertical sep=0.9cm,
        xticklabels at=edge bottom},
      width=8.4cm, height=3.6cm, scale only axis,
      view={0}{90},
      tick align=outside,
      tick label style={font=\footnotesize}, label style={font=\small},
      title style={font=\small},
      xmin=0, xmax=1, ymin=0, ymax=6.2832,
      ylabel={$\phi$},
      ytick={0,1.5708,3.1416,4.7124,6.2832},
      yticklabels={$0$,$\tfrac{\pi}{2}$,$\pi$,$\tfrac{3\pi}{2}$,$2\pi$},
      enlargelimits=false, grid=none,
      colormap/viridis,
      colorbar, colorbar style={font=\footnotesize, width=0.22cm,
        ylabel={$P(\phi\,|\,x)$}},
      point meta min=0, point meta max=1.2,
    ]
    \nextgroupplot[title={$\delta=|\Delta\alpha|E/p=1/3$（揃う）}]
      \addplot3[surf, shader=interp]
        table[x=x, y=phi, z=p]{\figdata/xphi_m6_b3_d0333333.dat};
      \draw[white, densely dashed, line width=0.7pt]
        (axis cs:0.81,0) -- (axis cs:0.81,6.2832);
      \draw[white, dotted, line width=0.7pt]
        (axis cs:0,1.5708) -- (axis cs:1,1.5708);
    \nextgroupplot[title={$\delta=10$（傾く）},
      xlabel={$x$ (1 周期に畳んだ位置)}]
      \addplot3[surf, shader=interp]
        table[x=x, y=phi, z=p]{\figdata/xphi_m6_b3_d10.dat};
      \draw[white, densely dashed, line width=0.7pt]
        (axis cs:0.81,0) -- (axis cs:0.81,6.2832);
      \draw[white, dotted, line width=0.7pt]
        (axis cs:0,1.5708) -- (axis cs:1,1.5708);
  \end{groupplot}
\end{tikzpicture}
\end{center}
\figexplain{$\frac{|\Delta \alpha|E}{p}=\frac{1}{3}$ では多くの場所で電場方向に揃うが、狭窄部では $x$ 軸に平行な方向へ棒の傾きが揃いやすい。一方、$\frac{|\Delta \alpha|E}{p}=10$ では全体として $x$ 軸方向に傾いており、狭窄部で向きを大きく変えずに通過できることを説明できる。}

\end{document}
